\subsection{Implizite Funktionen und Untermannigfaltigkeiten}

\begin{satz}{}
Es seien $ U \subset \mathbb{R}^n $, $ V \subset \mathbb{R}^m $ offene Mengen und $ f: U \times V \to \mathbb{R}^m $ eine $ C^1 $-Abbildung. Falls für ein Punkt $ (x_0, y_0) \in U \times V $ gilt:
\begin{enumerate}
    \item $ f(x_0, y_0) = 0 $,
    \item Die Jacobi-Matrix $ D_y f(x_0, y_0) $ ist invertierbar,
\end{enumerate}
dann existieren offene Mengen $ U' \subset U $ und $ V' \subset V $ mit $ (x_0, y_0) \in U' \times V' $ sowie eine differenzierbare Funktion $ g: U' \to V' $ mit


$$
\{(x,y) \in U' \times V' \mid f(x,y) = 0 \} = \{(x, g(x)) \mid x \in U' \}.
$$


\end{satz}

\begin{bemerkungbox}
Die Existenz von $ g $ basiert darauf, dass die Gleichung $ f(x,y) = 0 $ lokal eindeutig nach $ y $ aufgelöst werden kann, sofern die Ableitung $ D_y f(x,y) $ regulär ist.
\end{bemerkungbox}

\begin{korollar}
Falls $ f(x,y) $ eine $ C^r $-Funktion ist, dann ist auch $ g(x) $ eine $ C^r $-Funktion.
\end{korollar}

\begin{beweisbox}
Die Glattheit von $ g $ folgt aus der Anwendung der Kettenregel auf $ f(x,g(x)) = 0 $. Differenzieren nach $ x $ liefert


$$
D_x f(x, g(x)) + D_y f(x, g(x)) D_x g(x) = 0.
$$


Da $ D_y f(x, g(x)) $ invertierbar ist, kann nach $ D_x g(x) $ aufgelöst werden.
\end{beweisbox}

\begin{beispielbox}
\textbf{Beispiel:}  
Betrachten wir die Gleichung


$$
x^2 + y^2 - 1 = 0.
$$


Hier ist $ f(x,y) = x^2 + y^2 - 1 $. Da die partielle Ableitung


$$
D_y f(x,y) = 2y
$$


für $ y \neq 0 $ invertierbar ist, folgt, dass $ y $ lokal als Funktion von $ x $ geschrieben werden kann: 


$$
y = \pm \sqrt{1 - x^2}.
$$


\end{beispielbox}
